\section{Verdikt}
\label{sec:verdikt-azn}

Dieser Teil ist rot gesetzt. Rot hei{\ss}t: eine Empfehlung unter den in
Abschnitt~\ref{sec:annahmen-azn} erkl\"arten Annahmen, kein Befund. Wer nur
die schwarzen Teile dieses Dokuments liest, hat die belegte Lage ohne jede
Wertung.

\subsection{Votum}

\vd{Beobachten, nicht kaufen zum Kurs von \USDje{\AznKurs}. Die Schwelle
liegt bei \USDje{\AznSchwelleAmbition} auf einen Horizont von zehn Jahren;
das ist der Preis, zu dem die H\"urde von \Pct{\Huerde} erreicht wird, wenn
der freie Mittelzufluss mit der Rate w\"achst, die die eigene Umsatzambition
des Unternehmens f\"ur 2030 verlangt, und das Bewertungsvielfache konstant
bleibt.}

\vd{F\"ur den Halter: halten. Das Gesch\"aft ist intakt, die Dividende ist
gedeckt, und die Differenz zwischen Kurs und Schwelle betr\"agt
\Pct{\AznAbstandSchwelle}, zu wenig, um einen Verkauf mit Steuerfolge zu
rechtfertigen. F\"ur den Nichthalter: warten. Ein Aufbau bei
\USDje{\AznSchwelleAmbition} oder darunter ist eine andere Entscheidung als
ein Aufbau heute, und der Unterschied ist die gesamte Sicherheitsmarge.}

\subsection{Begr\"undung aus den Befunden}

\vd{Keine neue Zahl. Alles Folgende steht in den
Abschnitten~\ref{sec:abschluss-azn} bis \ref{sec:urteil-azn}.}

\vd{\emph{Was dagegen spricht, heute zu kaufen.} Das Basisjahr ist das beste
der sieben erfassten Jahre; jede Rechnung dieses Teils schreibt einen
Spitzenwert fort. Das gr\"o{\ss}te Einzelmedikament, \Pct{\AznFarxigaAnteil}
des Umsatzes, verliert gerade seinen Schutz, und wie viel Umsatz insgesamt an
auslaufenden Patenten h\"angt, ver\"offentlicht das Unternehmen nicht in den
hier vorliegenden Quellen. Der Kurs liegt \Pct{\AznAbstandSchwelle} \"uber der
Schwelle.}

\vd{\emph{Was daf\"ur spricht.} Eine Eigenkapitalrendite von
\Pct{\AznEigenkapitalrendite}, eine Verschuldung vom
\Faktor{\AznVerschuldungsgrad}-fachen des EBITDA, vier Jahre Prognosetreue
ohne Unterschreitung, und ein freier Mittelzufluss je Aktie, der zwischen 2019
und 2025 um \Pct{\AznCagr} pro Jahr gewachsen ist, dreimal so schnell wie
die \Pct{\AznWachstumVielfachZehn}, die der heutige Preis verlangt. Der Titel
hat den Branchenindex \"uber zehn Jahre geschlagen: \Pct{\AznRenditeZehn}
reine Kursrendite gegen \Pct{\Huerde} Gesamtrendite des Index. Die Dividende
kommt auf der Seite des Titels noch hinzu und ist in den
\Pct{\AznRenditeZehn} nicht enthalten.}

\vd{\emph{Was f\"urs Warten spricht statt f\"urs Weggehen.} Der Kurs steht
\PctBetrag{\AznAbstandHoch} unter dem Februarhoch, und was ihn gedr\"uckt hat,
ist ein angek\"undigtes Ereignis, das nun eingetreten ist. Ein weiterer
R\"uckgang um weniger als ein Zehntel bringt den Kurs an die Schwelle. Das ist
eine Entfernung, die ein einzelnes Quartal zur\"ucklegt.}

\vd{\emph{Wo dieses Votum vom Konsens abweicht.} Es weicht in der Richtung
ab, nicht nur im Niveau. \Stueck{\AznKonsensHaeuser} H\"auser sagen
\enquote{Strong Buy}, dieses Dokument sagt beobachten; das mittlere Kursziel
von \USDje{\AznKonsensZiel} liegt \Pct{\AznKonsensUeberSchwelle} \"uber der
strengsten und \Pct{\AznKonsensUeberVariante} \"uber der mittleren Schwelle
dieses Berichts. Wenn eine ganze Abdeckung \"uber der eigenen Schwelle liegt,
ist die wahrscheinlichste Erkl\"arung eine eigene Annahme und nicht ein
kollektiver Irrtum, und sie ist hier benennbar: Ein Kursziel auf zw\"olf
Monate bewertet mit einem Gewinnvielfachen und unterstellt damit einen
Fortf\"uhrungswert, den eine Rechnung mit Buchwert als Endwert nicht
enth\"alt. Der Unterschied ist die Bewertungspr\"amie, nicht die
Ertragslage. Deshalb steht in Tabelle~\ref{tab:azn-schwelle-variante} eine
zweite und in Abschnitt~\ref{sec:schwelle-azn} eine dritte Endwertannahme:
Erst die dritte, die das Vielfache konstant h\"alt, liegt in derselben
Gr\"o{\ss}enordnung wie der Markt, und sie ergibt die Schwelle, die das
Votum tr\"agt.}

\subsection{Einstiegsschwelle}
\label{sec:schwelle-azn}

\vd{Drei Endwertannahmen, drei Schwellen auf zehn Jahre. Sie stehen
nebeneinander, weil die Wahl unter ihnen das Ergebnis bestimmt und der Leser
sie selbst treffen k\"onnen soll.}

\begin{table}[htbp]
\centering
\caption{Einstiegsschwelle auf zehn Jahre bei einer H\"urde von
\Pct{\Huerde}, je Endwertannahme. Kurs heute: \USDje{\AznKurs}.}
\label{tab:azn-schwellen-drei}
\begin{tabular}{llr}
\toprule
Endwertannahme & unterstelltes Wachstum & Schwelle\\
\midrule
Buchwert je Aktie            & keines & \USDje{\AznSchwelleZehn}\\
Kurs von heute, nominal      & keines & \USDje{\AznSchwelleVarZehn}\\
konstantes Vielfaches        & Umsatzambition 2030 & \USDje{\AznSchwelleAmbition}\\
\bottomrule
\end{tabular}
\end{table}

\vd{Die dritte Zeile tr\"agt das Votum, und sie ist zugleich die
anspruchsvollste: Sie unterstellt, dass das Unternehmen die eigene Ambition
erreicht (\Pct{\AznAmbitionRate} Umsatzwachstum pro Jahr bis 2030), dass der
freie Mittelzufluss im Gleichschritt mitw\"achst und dass der Markt in zehn
Jahren dasselbe Vielfache zahlt wie heute. Unter diesen drei wohlwollenden
Annahmen ist der faire Preis \USDje{\AznSchwelleAmbition}, und der Kurs liegt
\Pct{\AznAbstandSchwelle} dar\"uber.}

\begin{table}[htbp]
\centering
\caption{Gefordertes Wachstum des Mittelzuflusses in Prozent pro Jahr bei
konstantem Bewertungsvielfachen. Die Rate h\"angt nicht vom Horizont ab;
das ist eine Eigenschaft der Rechnung, kein Rundungseffekt.}
\label{tab:azn-vielfach}
\begin{tabular}{lrrr}
\toprule
Niveau & 5 Jahre & 10 Jahre & 15 Jahre\\
\midrule
\tabellenkoerper{data/flow_azn_vielfach}
\bottomrule
\end{tabular}
\end{table}

\vd{Zum Vergleich: erreicht wurden \Pct{\AznCagr} pro Jahr zwischen 2019 und
2025. Die geforderte Rate ist also gut erreichbar. Das macht den Kurs nicht
billig, sondern erkl\"art, warum er kein Schn\"appchen ist: Das Wachstum ist
eingepreist, nicht mehr und nicht weniger.}

\FloatBarrier

\subsection{Ausl\"oser, die das Votum drehen}

\vd{Alle vier sind Gr\"o{\ss}en, die das Unternehmen selbst ver\"offentlicht
oder die auf dem Kurszettel stehen. Ein Ausl\"oser, den niemand nachsieht,
ist Dekoration.}

\vd{\emph{Zum Kauf.} Ein Kurs unter \USDje{\AznSchwelleAmbition}. Oder ein
freier Mittelzufluss des Gesch\"aftsjahres 2026 \"uber dem Niveau von 2025
trotz des Patentablaufs bei Farxiga; das w\"urde belegen, dass die Pipeline
den Ausfall tr\"agt, und die Schwelle hebt sich entsprechend.}

\vd{\emph{Zum Verkauf.} Eine Senkung der Prognose f\"ur 2026 unterhalb des
mittleren einstelligen Umsatzwachstums; nach vier Jahren ohne Unterschreitung
w\"are das ein Bruch mit dem Muster und kein Rauschen. Oder ein Wegfall der
Ausnahme von den Z\"ollen nach \ZollParagraf{}, die im Oktober 2025 f\"ur drei
Jahre vereinbart wurde.}

\subsection{Was dieses Votum nicht wissen kann}

\vd{Die Gr\"o{\ss}e, an der der Fall h\"angt, ist nicht berechenbar: wie viel
des heutigen Umsatzes in den n\"achsten f\"unf Jahren durch Patentabl\"aufe
wegf\"allt und wie viel die Pipeline davon ersetzt. Das Unternehmen
ver\"offentlicht die Aufstellung gesondert, sie liegt hier nicht vor, und eine
Sch\"atzung an ihrer Stelle w\"are genau die Sorte plausibel aussehender Zahl,
die dieser Bericht nicht produziert.}

\vd{Drei Argumente gegen das eigene Votum, ausdr\"ucklich:}

\vd{Erstens h\"angt die Schwelle an der Annahme eines konstanten
Bewertungsvielfachen. Ein Unternehmen, das seine Ambition erreicht und
zwanzig neue Medikamente einf\"uhrt, k\"onnte in zehn Jahren ein h\"oheres
Vielfache tragen; dann ist die Schwelle zu niedrig angesetzt und das Votum zu
streng.}

\vd{Zweitens ist die Reihe, aus der die \Pct{\AznCagr} stammen, nur sieben
Jahre lang und beginnt in einem schwachen Jahr. Eine l\"angere Reihe w\"urde
eine niedrigere Rate ergeben und das Argument \enquote{das Unternehmen kann
das Geforderte leisten} schw\"achen.}

\vd{Drittens sagen \Stueck{\AznKonsensHaeuser} H\"auser \enquote{Strong Buy}.
Ein Votum, das eine einhellige Abdeckung ignoriert, ist nicht verteidigbar;
eines, das vor ihr einknickt, war nie eines. Dieses hier bleibt bei
beobachten, und der Grund steht oben: Es misst gegen die H\"urde einer
Branchenrendite \"uber zehn Jahre, nicht gegen ein Kursziel \"uber zw\"olf
Monate. Wer den k\"urzeren Ma{\ss}stab anlegt, kommt legitim zu einem anderen
Ergebnis.}
